% !PRINTED 82
\weberpage{82}
%<<CONT
oder durch Integration:
\[
\vartheta_{11}' = c\,\vartheta_{00}\vartheta_{10}\vartheta_{01},
\]
worin $c$ von $\omega$ unabhängig ist. Durch \S\,20, 21 waren die $\vartheta$-Funktionen bestimmt bis auf einen von $v$, $\omega$ unabhängigen, allen gemeinschaftlichen Faktor. \gesperrt{Über diesen Faktor soll nun so verfügt werden, daß die Konstante $c$ den Wert $\pi$ erhält,} also die Formel besteht:
\weq{5}{\vartheta_{11}' = \pi\,\vartheta_{00}\vartheta_{10}\vartheta_{01},}
\gesperrt{und dadurch sind jetzt die $\vartheta$-Funktionen bis auf das gemeinschaftliche Vorzeichen $\pm$ vollständig definiert.} Durch Anwendung von (5) läßt sich der Formel (1) die Gestalt geben:
\weq{6}{\vartheta_{11}'(v)\vartheta_{01}(v) - \vartheta_{01}'(v)\vartheta_{11}(v) = \pi\,\vartheta_{01}^2\,\vartheta_{10}(v)\vartheta_{00}(v),}
und ebenso erhält man mit Benutzung von \S\,21, (8):
\weq{7}{\vartheta_{11}'(v)\vartheta_{00}(v) - \vartheta_{00}'(v)\vartheta_{11}(v) = \pi\,\vartheta_{00}^2\,\vartheta_{10}(v)\vartheta_{01}(v),}
\weq{8}{\vartheta_{11}'(v)\vartheta_{10}(v) - \vartheta_{10}'(v)\vartheta_{11}(v) = \pi\,\vartheta_{10}^2\,\vartheta_{01}(v)\vartheta_{00}(v),}
\weq{9}{\vartheta_{10}'(v)\vartheta_{01}(v) - \vartheta_{01}'(v)\vartheta_{10}(v) = -\pi\,\vartheta_{00}^2\,\vartheta_{00}(v)\vartheta_{11}(v),}
\weq{10}{\vartheta_{00}'(v)\vartheta_{01}(v) - \vartheta_{01}'(v)\vartheta_{00}(v) = -\pi\,\vartheta_{10}^2\,\vartheta_{10}(v)\vartheta_{11}(v),}
\weq{11}{\vartheta_{00}'(v)\vartheta_{10}(v) - \vartheta_{10}'(v)\vartheta_{00}(v) = \pi\,\vartheta_{01}^2\,\vartheta_{01}(v)\vartheta_{11}(v).}

Differentiieren wir die Gleichungen (9), (10), (11) nach $v$ und setzen $v=0$, so erhalten wir mittels (4) und (5) die Relationen
\weq{12}{4\,\frac{d}{d\omega}\log\frac{\vartheta_{10}}{\vartheta_{01}} = i\pi\,\vartheta_{00}^4,}
\weq{13}{4\,\frac{d}{d\omega}\log\frac{\vartheta_{00}}{\vartheta_{01}} = i\pi\,\vartheta_{10}^4,}
\weq{14}{4\,\frac{d}{d\omega}\log\frac{\vartheta_{10}}{\vartheta_{00}} = i\pi\,\vartheta_{01}^4.}

Wenn man die Fundamentformeln % [sic] printed \textquotedbl{}Fundamentformeln\textquotedbl{}, expected Fundamentalformeln
für die $\vartheta$-Funktionen zweimal logarithmisch differentiiert, so erkennt man, daß die Funktionen
\[
\vartheta^2(v)\,\frac{d^2\log\vartheta(v)}{dv^2} = \vartheta''(v)\vartheta(v) - \vartheta'(v)\vartheta'(v)
\]
$\Theta$-Funktionen zweiter Ordnung mit der Charakteristik $(0,0)$ sind, und man kann sie daher durch die Funktionen $\vartheta^2(v)$ linear ausdrücken. Auf diese Weise ergibt sich
%CONT>>
